%----------------------------------------------------------------------------------------
%	PACKAGES AND OTHER DOCUMENT CONFIGURATIONS
%----------------------------------------------------------------------------------------
\documentclass[a4paper,12pt]{article}
\usepackage[english]{babel}
\usepackage[latin1]{inputenc}
\usepackage{amsmath}
\usepackage{xcolor}
\usepackage{amssymb}
\usepackage{amsfonts}
\usepackage{graphicx}
\usepackage{geometry}
\definecolor{link}{rgb}{0.4, 0.2, 0.2}
\geometry{
	a4paper,
	total={170mm,257mm},
	left=20mm,
	top=20mm,
}
\usepackage{hyperref}
\hypersetup{
	unicode=false,          % non-Latin characters in Acrobat’s bookmarks
	pdftoolbar=true,        % show Acrobat’s toolbar?
	pdfmenubar=true,        % show Acrobat’s menu?
	pdffitwindow=false,     % window fit to page when opened
	pdfstartview={FitH},    % fits the width of the page to the window
	pdftitle={Research Planning Report},    % title
	pdfauthor={Aaron Panaitescu},     % author
	pdfkeywords={neuroplasticity, Parkinson's, stimulation}, % list of keywords
	pdfnewwindow=true,      % links in new PDF window
	colorlinks=true,       % false: boxed links; true: colored links
	linkcolor=link,          % color of internal links (change box color with linkbordercolor)
	citecolor=link,        % color of links to bibliography
	filecolor=link,      % color of file links
	urlcolor=link           % color of external links
}

\usepackage{xspace} 

\newcommand{\ballarrow}{\tikz[baseline=-0.6ex]{\draw[->] (0,0) -- (0.5,0); \filldraw (0.5,0) circle (2pt);}}
\newcommand{\Na}{$Na^+$\xspace}
\newcommand{\K}{$K^+$\xspace}
\newcommand{\Ca}{$Ca^{2+}$\xspace}
\hyphenpenalty=750


\begin{document}
\begin{titlepage}
	\newcommand{\HRule}{\rule{\linewidth}{0.5mm}}
	\setlength{\topmargin}{0in}
	\center

	\begin{flushleft} \large
		\begin{minipage}{0.4\textwidth}
			\includegraphics[scale=0.14]{imperial.png}
		\end{minipage}
	\end{flushleft}
	\vspace{4cm}

	%----------------------------------------------------------------------------------------
	%	TITLE SECTION
	%----------------------------------------------------------------------------------------
	\textbf{\large Literature review and thesis proposal}\\[0.1cm]
	{\large MRes. Neurotechonlogy}\\[0.5cm]

	\HRule \\[0.4cm]
	{\Large \bfseries Investigating plasticity in \\ Cortico-Basal Ganglia-Thalamus models \\ to improve stimulation-based treatments }
	\HRule \\[1cm]

	%----------------------------------------------------------------------------------------
	%	AUTHOR SECTION
	%----------------------------------------------------------------------------------------

	{\large Aaron Panaitescu \\
	(\textit{CID:} 02054726) \\[0.4cm]
	\textbf{Supervisor:} \\
	Dr. Hayriye Cagnan}

	%----------------------------------------------------------------------------------------
	%	DATE SECTION
	%----------------------------------------------------------------------------------------

	\vfill
	{Word Count: ????}\\[0.4cm]
	{\large \today}\\[0.8cm]
\end{titlepage}
%----------------------------------------------------------------------------------------
%	REPORT CONTENT
%----------------------------------------------------------------------------------------
\tableofcontents
\newpage

\section{Introduction / Abstract?}
**Stylistically, I would prefere to do an abstract here and breakdown the concepts in the
following section**
\section{Background and literature review}
**start this off nicely with a diagram of the research gap**

\subsection{Parkinson's Disease}
Parkinson's Disease (PD) is characterized by the loss of dopaminergic neurons and the presence of
\emph{Lewy Bodies} (LB) in the \emph{Substantia Nigra pars Compacta} (SNc)
\cite{del2018advances}. This dopaminergic depletion in the SNc leads to striatal dopamine
deficiency, which in turn leads to excessive synchronization in the basal ganglia.
Hyper-synchrony is a hallmark of the PD \cite{hammond2007pathological, helmich2012cerebral},
and thus, can be a promising therapeutic target.

\subsection{Deep brain stimulation: theory and practice}
Tremor is primarly generated in the Cortical-Basal Ganglia-Thalamic (CBGT) network.
Thus, deep brain stimulation (DBS) in this region emerged as an effective tool in
the treatment of advanced PD patients \cite{del2018advances}.
DBS works through the implantation of electrodes in particular regions
of the brain.
The intuition behind this method is that it enables stimulation driven decorelatation of the
different basal ganglia nuclei, pushing it to a non-pathological state.
There is significant evidence that performing DBS in the STN or the GPi, can have major quality
of life improvements in PD patients, reducing tremor and restoring motor control
\cite{rodriguez2005bilateral}.
DBS does, however have significant limitations.

Firstly, like any highly-invasive procedure, the implatation procedure has many risks.
Advances in electrode technology and surgical techniques have reduced some of these risks.
*cite*

Secondly, patient symptoms are only ameloriated while the stimulation is on. This means that
the device needs to be on continously, which can induce numerous side effects in patients
**citeee** and also creates the need for a continous power supply.

Finally, the mechanism through which DBS reduces tremors is not fully understood yet, there
are many questions left to answer. For instance, why does DBS in PD get effective results around
130 Hz, when tremor frequency is on average around 20 Hz **check this**.

Therfore, several areas of research emerge in improving stimulation-based treatments of PD:
finding non-invasive alternatives \cite{saturnino2017target, schwab2020spike}, eleciting lasting,
plastic changes in CBGT circuit, closing the loop and coupling stimulation to the presence of
patient symptoms \cite{beudel2018adaptive}, and improving stimulation patterns to harness the
exploit the rythmic patterns activity in the basal ganglia \cite{cagnan2017stimulating, west2022stimulating}.

\subsection{Stimulating at the right time}
The Parkinsonian state is modulated by the interplay between four pathways in the CBGT
the \textbf{hyperdirect} excitatory cortico-subthalamic pathway,
the \textbf{direct} inhibitory striato-pallidal pathway,
the \textbf{indirect} inhibitory subthalamic-pallido pathway and
the inhibitory \textbf{pallido-subthalamic} (PS) pathway. In the Parkinsonian state the hyperdirect
pathway is downregulated, while pallido-subthalamic pathway is upregulated, leading to a shift in
beta band oscillatory activity from sub-band $\beta_1$ to sub-band $\beta_2$, which corellates to
hyper-synchrony in the basal ganglia \cite{west2022stimulating}.
** \textbf{add diagram} **

**important** \cite{cagnan2017stimulating} \cite{beudel2018adaptive} \cite{west2022stimulating}

\subsection{Plasticity to recover network states}
Neuroplasticity can be defined as "the ability of the nervous system to change its activity in
response to intrinsic or extrinsic stimuli by reorganizing its structure, functions, or
connections" \cite{mateos2019impact}.
Plasticity occurs at different scales in the brain, ranging from molecular to neural networks.
The focus of this project is synaptic plasticity, which occurs between a pre-synaptic and a
pos-synaptic neuron. While, the mechanisms of plasticity have been throughly studied **cite**,
this understanding is just starting to find its way into clinical applications in
neurorehabilitation (e.g. stroke, trauma, spinal cord injury \cite{cramer2011harnessing}) and
neuroprostheses \cite{lebedev2017brain}. \\
**should mention STDP above** \\
**OMG I'm cooked how tf did I miss \cite{hauptmann2009cumulative, hauptmann2010restoration}?!** \\
**\textbf{coordinated reset (CR) stimulation}** \\
**this should be relevant as well** \cite{schwab2020spike} \\

\subsection{Neuron-level vs. Mean-field models}
Neuronal dynamics models are defined by defined by rules (ODEs) of how voltage in an unit changes
over time in response to voltage in the network and external current (e.g. noise, stimulation
current).
The question then becomes what is an unit? There are two broad categories branching out of this
question: An unit can be an individual neuron or a cluster of neurons. This section introduces
some common types of models, justifying the direction chosen for this project.

\paragraph{Neuron models} (or spiking models) work by combining networks of units that individually
simulate neuronal action potentials (APs).
These networks can then be organized into different nuclei (e.g. STN) and coherced to respect
excitatory or inhibitory connections between nuclei.
Within this class of models, different schemes for modeling neuronal dynamics exist,
such as Integrate-And-Fire (IF) \cite{gerstner2014if} and Hodgkin-Huxley
\cite{hodgkin1952measurement, gerstner2014hh}.

\paragraph{Mean-field models} simulate the local field potentials (LFPs) of clusters of neurons
and connect these clusters with population-level connectivity rules. (e.g.
\cite{jansen1995electroencephalogram})

% Neuron-level
\subsubsection{Integrate-And-Fire}
IF models treat action potentials as events.
This reduction is justified by the fact that the shape of APs is always \textit{approximately}
the same, meaning they \textit{cannot} convey information. \cite{gerstner2014if} describe a leaky IF model
is represented by an electrical circuit with a resistor($R$) and a capacitor($C$) in parallel
($I(t) = I_R + I_C$), combined with a reset condition when the potential exceeds threshold
$\theta$.
This can be modeled by,
\begin{align}
	\tau_m \dot u                                           & = -[u(t) - u_{rest}] + RI(t),                                    \\
	\intertext{Together with the reset condition,}
	\lim_{\delta \rightarrow 0; \delta > 0} u(t^f + \delta) & = u_r,                                                           \\
	t^f                                                     & = \{t | u(t)                                        = \theta \}.
\end{align}
Where $u$ is the membrane voltage, $u_{rest}$ is the resting potential, $u_r$ is the reset
potential, $\tau_m = RC$ is the time constant, and $I$ is external current.

These equations describe the update rules in a single unit. Connectivity rules between units
can be added to form neural networks. These connections can then be strengthend or weakend
depending on spike timings between units, modeling STDP. \cite{shupe2021integrate} simulate
plasticity in cortical columns; \cite{kromer2023synaptic} model plasticity between basal ganglia
nuclei both using a IF model.

\subsubsection{Hodgkin-Huxley}
The Hodgkin-Huxley (HH) model \cite{hodgkin1952measurement} simulates membrane potential by
considering the dynamics of different ion conductances.
Each ion has a reversal potential at which the force of the electrical and chemical gradients
across the membrane even out.
These can be modeled in an electrical circuit as a battery and the total channels of a
particular ion can be modeled as a resistor. Each modelled ion has a parallel branch in the
corresponding electrical circuit. The following equations adapted from \cite{gerstner2014hh}
describe the HH model:
\begin{align}
	I(t) & = I_C(t) + \sum_k I_k(t),
	\intertext{Where the current of each ion $k$ can be described in terms of the open probability
		of a $k$ channel, $p_k$ and the maximum conductance when all $k$ channels are open, $g_k$:}
	I_k  & = g_k p_k(u - E_k).
\end{align}
In particular Hodgkin and Huxley modeled $Na^+$, $K^+$ and a leak $L$ currents:
\begin{align}
	C \dot u = g_{Na}m^3h(E_{Na} - u) + g_K n^4(E_K - u) + g_L (E_L - u),
\end{align}
Where $m$, $h$ and $n$ are open probabilities for subunints of $Na$ or $K$ channels governed by:
\begin{align}
	\dot x = - \frac 1 {\tau_m(u)}[x - x_0(u)], &  & x \in \{m, h, n\}.
\end{align}
As in the IF model these rules describe the dynamics of a single unit, which can subsequently
be connected to model a networks of neurons (e.g. \cite{terman2002activity}) and augumented to
include plasticity rules \cite{borges2016effects}.

\subsubsection{Rubin-Terman}
\cite{terman2002activity} augumented the HH model to create a subthalamopallido network of the
basal ganglia (RT model), where GPe and STN neurons are modeled according following equation,
\begin{align}
	C \dot u = - I_L - I_K - I_{Na} - I_T - I_{Ca} - I_{AHP} - I_{s \rightarrow d} -
	I_{s \rightarrow s} + I_{app}
\end{align}
Where leak ($L$), \K, \Na and high-threshold \Ca currents \cite{song2000characterization}, are
described in HH-like equations, and low threshold calcium current $I_T$ is different for GPe and
STN neurons. $I_{AHP}$ represents \Ca-activated, voltage-independent \textit{afterhyperpolarization}
(AHP) $K^+$ current. $I_{s \rightarrow d}$ represents the inter-populational influence and is
defined by
$I_{s \rightarrow d} = \sum_{j} I_{ij}^{ds}$ where $s, d \in \{\text{GPe}, \text{STN}\}$.
$I_{ij}^{ds}$ is the current from the presynaptic neuron $j$ in $s$ to the
postsynaptic neuron $i$ in $d$, and is defined by
$I_{ij}^{ds} = w_{ij}^{ds}s_{ij}^{ds}(t - \tau_{ds})(u_i - E_{ds})$ where $w_{ij}$ is the
synaptic strength, $s_{ij}^{ds}$ is a synaptic variable and $\tau_{ds}$ is the transimission
delay between $d$ and $s$.
\cite{madadi2022inhibitory} expand on this model by adding inhibitory spike-timing-dependent
plasticity to this network with $\Delta w = \eta(\exp(-|\Delta t| / \tau_w) - \alpha)$,
where $\eta$ is the \textit{learning rate}, $\tau_w$ is the time constant of the plasticity
rule exponential decay and $\alpha$ is the depression factor. \cite{rubin2012basal} expand
the RT model into a CBGT model and analyze ... \\
**\cite{hauptmann2009cumulative, hauptmann2010restoration} $\rightarrow$ GPe-STN + plasticity +
CR stimulation** \\
**\cite{rubin2012basal} (on multi-site CR stimulation) \emph{"in its current formulation, does
	depend critically on spike timing dependent plasticity inconnections among excitatory neurons
	within the STN, the presence of which currently lacks experimental confirmation. It is possible
	that the general theory is correct but that the site of plasticity is elsewherein the BG, an
	idea that remains to be explored"}**

% Population
\subsubsection{Jansen}
**this should be brief** \cite{jansen1995electroencephalogram} \cite{west2022stimulating} \\
**phase-dependent plasticity** \cite{duchet2023mean} \\


\subsection{**Other ways of improving stimulation-based treatments**}
**stimulation parameter optimizations, closing the loop (e.g. aDBS \cite{beudel2018adaptive})**
\textit{
	Maybe this can be folded into the stimulating at the right time part
	since they are pretty closely related. Also should mention the idea that
	these methods are not mutually exclusive, meaning that in principle they
	could be combined, adaptive stim closing the loop, with inducing plastic
	changes as target
}

\section{Project Plan}


\subsection{Aims}
\begin{enumerate}
	\item Model neuroplasticity in a Parkinsonian CGBT network
	\item Investigate the viabillity of harnessing plasticity to remove the
	      system from the pathological state and analyze the dynamics that follow

	      **here i care about things like for how long and how does the network change.
	      To what degree can we induce changes etc.**

	      **should look into viable timescales + noise analyze how long does it take to rebound to
	      pathological state**
	\item Try to link potential results to potential stimulation protocols?
\end{enumerate}

\subsection{Methodolgy} **How indeed?** HH/IF Pakrkinsoni model + plasticity rules, trying
different stimulation-based protocols (link with experimental data?)

**look into how to tune IF model parameters from experimental data**

\subsection{Timeline}
**Gant chart thingy**



%**********************************************%
\newpage
%----------------------------------------------------------------------------------------
%	REFERENCES
%----------------------------------------------------------------------------------------
\addcontentsline{toc}{section}{References}
\bibliography{refs}
\bibliographystyle{apalike}
\newpage

%----------------------------------------------------------------------------------------
%	APPENDIX  
%----------------------------------------------------------------------------------------
\appendix
\section*{Appendices}
\addcontentsline{toc}{section}{Appendices}
\renewcommand{\thesubsection}{\Alph{subsection}}
\subsection{First appendix}
\subsection{Second appendix}

\end{document}
